% answers.tex —— 1977 年·江苏南通市初试卷 · 参考答案与详细解析片段
% 由 tools/generate_projects.py 依来源总账自 1977/1977_数学真题_江苏南通市初试_试题.md + 1977/1977_数学真题_江苏南通市初试_解析正文.md + 1977/1977_数学真题_江苏南通市初试_答案要点.md 生成，勿手改（改源 MD 后重跑）


\begin{gktip}
固定源 PDF 仅含 3 页题面，无解答可对照。以下各题均由编者独立推导，最终结论与答案要点逐题一致；原卷大题的中文数字编号已转写为阿拉伯数字，各大题下的小题编号「1．2．3．」已转写为（1）（2）（3）。第 3 题的（1）（3）两小问按原卷题面推导均无实数解（（1）整理后判别式小于零，（3）候选根不满足对数底数条件），系按原卷如实推得，非转录疑误。
\end{gktip}

\begin{answer}
\gksolhead{1}{}

\textbf{题目}：

化简、计算：

（1） $\dfrac{-2mn}{m-3n}\sqrt{\dfrac{(m+3n)^2}{2mn}-6}$；$(0<m<3n)$

（2） $\left(\dfrac14\right)^{-1.5}\left[\left(\dfrac{7}{11}\right)^0\right]^{-3}-\left(\dfrac35\right)^{-2}\div\left(2\dfrac79\right)^{\frac12}$。

\Needspace{4\baselineskip}
\textbf{解析}

\gksubpart{（1）}

\Needspace{5\baselineskip}
由 $0<m<3n$ 知 $m-3n<0$、$2mn>0$。先化简根号内：

\nopagebreak[4]
\[
\frac{(m+3n)^2}{2mn}-6
=\frac{m^2+6mn+9n^2-12mn}{2mn}
=\frac{(m-3n)^2}{2mn}.
\]

\Needspace{5\baselineskip}
因 $m-3n<0$，故

\nopagebreak[4]
\[
\sqrt{\frac{(m-3n)^2}{2mn}}=\frac{3n-m}{\sqrt{2mn}}.
\]

\Needspace{5\baselineskip}
于是

\nopagebreak[4]
\[
\frac{-2mn}{m-3n}\cdot\frac{3n-m}{\sqrt{2mn}}
=\frac{-2mn}{-(3n-m)}\cdot\frac{3n-m}{\sqrt{2mn}}
=\frac{2mn}{\sqrt{2mn}}
=\sqrt{2mn}.
\]

\gksubpart{（2）}

\Needspace{5\baselineskip}
逐项计算：

\nopagebreak[4]
\[
\left(\frac14\right)^{-1.5}=4^{1.5}=8,\qquad
\left[\left(\frac{7}{11}\right)^0\right]^{-3}=1^{-3}=1,
\]

\nopagebreak[4]
\[
\left(\frac35\right)^{-2}=\frac{25}{9},\qquad
\left(2\frac79\right)^{\frac12}=\left(\frac{25}{9}\right)^{\frac12}=\frac53.
\]

\Needspace{5\baselineskip}
所以原式

\nopagebreak[4]
\[
=8\times1-\frac{25}{9}\div\frac53
=8-\frac{25}{9}\times\frac35
=8-\frac53
=\frac{19}{3}.
\]

\gkanswer{答案（1）：$\sqrt{2mn}$。}

\gkanswer{答案（2）：$\dfrac{19}{3}$。}

\end{answer}
\begin{answer}
\gksolhead{2}{}

\textbf{题目}：

在实数范围内分解因式：

（1） $a^4-10a^2+16$；

（2） $x^2-4xy+4y^2-3x+6y+2$。

\Needspace{4\baselineskip}
\textbf{解析}

\gksubpart{（1）}

\Needspace{5\baselineskip}
把 $a^4-10a^2+16$ 看作关于 $a^2$ 的二次三项式：

\nopagebreak[4]
\[
a^4-10a^2+16=(a^2-2)(a^2-8).
\]

\Needspace{5\baselineskip}
在实数范围内继续分解：

\nopagebreak[4]
\[
(a^2-2)(a^2-8)=(a+\sqrt2)(a-\sqrt2)(a+2\sqrt2)(a-2\sqrt2).
\]

\gksubpart{（2）}

\Needspace{5\baselineskip}
前三项配成完全平方：

\nopagebreak[4]
\[
x^2-4xy+4y^2-3x+6y+2=(x-2y)^2-3(x-2y)+2.
\]

\Needspace{5\baselineskip}
再把 $x-2y$ 视为整体分解：

\nopagebreak[4]
\[
(x-2y)^2-3(x-2y)+2=(x-2y-1)(x-2y-2).
\]

\gkanswer{答案（1）：$(a+\sqrt2)(a-\sqrt2)(a+2\sqrt2)(a-2\sqrt2)$。}

\gkanswer{答案（2）：$(x-2y-1)(x-2y-2)$。}

\end{answer}
\begin{answer}
\gksolhead{3}{}

\textbf{题目}：

求 $x$：

（1） $\dfrac{2x+1}{x-3}-\dfrac{5x-11}{x^2-4x+3}=\dfrac{3}{x-1}$；

（2） $\sqrt{2x-4}-\sqrt{x+5}=1$；

（3） $\log_{(x-1)}(2x^2-3x-1)=2$。

\Needspace{4\baselineskip}
\textbf{解析}

\gksubpart{（1）}

\Needspace{5\baselineskip}
定义域要求 $x\ne1$、$x\ne3$。由 $x^2-4x+3=(x-1)(x-3)$，两边同乘 $(x-1)(x-3)$：

\nopagebreak[4]
\[
(2x+1)(x-1)-(5x-11)=3(x-3).
\]

\Needspace{5\baselineskip}
展开整理：

\nopagebreak[4]
\[
2x^2-x-1-5x+11=3x-9,
\qquad
2x^2-9x+19=0.
\]

\Needspace{5\baselineskip}
判别式

\nopagebreak[4]
\[
\Delta=(-9)^2-4\times2\times19=81-152=-71<0,
\]

故原方程无实数根。

\gksubpart{（2）}

\Needspace{5\baselineskip}
根号要求 $x\ge2$。设 $u=\sqrt{2x-4}$、$v=\sqrt{x+5}$（$u$、$v\ge0$），则 $u-v=1$，且

\nopagebreak[4]
\[
2v^2-u^2=2(x+5)-(2x-4)=14.
\]

\Needspace{5\baselineskip}
把 $u=v+1$ 代入：

\nopagebreak[4]
\[
2v^2-(v+1)^2=14
\quad\Longrightarrow\quad
v^2-2v-15=0
\quad\Longrightarrow\quad
v=5\ (v\ge0).
\]

由 $v=\sqrt{x+5}=5$ 得 $x=20$。检验：$\sqrt{2\times20-4}-\sqrt{20+5}=6-5=1$，成立。

\gksubpart{（3）}

\Needspace{5\baselineskip}
对数成立要求底数 $x-1>0$ 且 $x-1\ne1$，即 $x>1$ 且 $x\ne2$。由对数定义，

\nopagebreak[4]
\[
2x^2-3x-1=(x-1)^2
\quad\Longrightarrow\quad
x^2-x-2=0
\quad\Longrightarrow\quad
x=2\ \text{\text{或}}\ x=-1.
\]

$x=2$ 使底数为 $1$，$x=-1$ 使底数为 $-2$，均不满足底数条件，故原方程无解。

\gkanswer{答案（1）：无实数根（去分母整理得 $2x^2-9x+19=0$，$\Delta=-71<0$；$x\ne1$、$x\ne3$）。}

\gkanswer{答案（2）：$x=20$。}

\gkanswer{答案（3）：无解（候选根 $x=2$、$x=-1$ 均不满足对数底数条件）。}

\end{answer}
\begin{answer}
\gksolhead{4}{}

\textbf{题目}：

求下列函数的定义域：

（1） $y=\dfrac{1}{1-\sqrt{x-1}}$；

（2） $y=\operatorname{lg}\operatorname{tg}x$。

\Needspace{4\baselineskip}
\textbf{解析}

\gksubpart{（1）}

\Needspace{5\baselineskip}
需要 $x-1\ge0$，且分母 $1-\sqrt{x-1}\ne0$（即 $\sqrt{x-1}\ne1$，$x\ne2$）。故定义域为

\nopagebreak[4]
\[
x\ge1\ \text{\text{且}}\ x\ne2.
\]

\gksubpart{（2）}

\Needspace{5\baselineskip}
需要 $\operatorname{tg}x>0$，即 $x$ 位于第一或第三象限：

\nopagebreak[4]
\[
k\pi<x<k\pi+\frac{\pi}{2}\quad(k\ \text{\text{为整数}}).
\]

\gkanswer{答案（1）：$x\ge1$ 且 $x\ne2$。}

\gkanswer{答案（2）：$k\pi<x<k\pi+\dfrac{\pi}{2}$（$k$ 为整数）。}

\end{answer}
\begin{answer}
\gksolhead{5}{}

\textbf{题目}：

不查表：

（1） 求 $4^{1-\log_2\frac{2}{5}+\log_2\sqrt{2}}$，

（2） 解三角形：已知 $\triangle ABC$ 中，$\angle B=120^\circ$，$\angle C=15^\circ$，$b=100$，求 $c$（可用根式表示结果）。

（3） 将直角坐标方程 $x^2-y^2=a^2$ 化成极坐标方程。

\Needspace{4\baselineskip}
\textbf{解析}

\gksubpart{（1）}

\nopagebreak[4]
\[
4^{1-\log_2\frac25+\log_2\sqrt2}
=4\times4^{-\log_2\frac25}\times4^{\log_2\sqrt2}.
\]

\Needspace{5\baselineskip}
由 $4=2^2$，

\nopagebreak[4]
\[
4^{-\log_2\frac25}=2^{-2\log_2\frac25}=\left(\frac25\right)^{-2}=\frac{25}{4},
\qquad
4^{\log_2\sqrt2}=2^{2\log_2\sqrt2}=(\sqrt2)^2=2.
\]

故原式 $=4\times\dfrac{25}{4}\times2=50$。

\gksubpart{（2）}

\nopagebreak[4]
\[
\angle A=180^\circ-120^\circ-15^\circ=45^\circ.
\]

\Needspace{5\baselineskip}
由正弦定理，

\nopagebreak[4]
\[
c=\frac{b\sin C}{\sin B}
=\frac{100\sin15^\circ}{\sin120^\circ}.
\]

\Needspace{5\baselineskip}
代入 $\sin15^\circ=\dfrac{\sqrt6-\sqrt2}{4}$、$\sin120^\circ=\dfrac{\sqrt3}{2}$：

\nopagebreak[4]
\[
c=\frac{100\times\dfrac{\sqrt6-\sqrt2}{4}}{\dfrac{\sqrt3}{2}}
=50\times\frac{\sqrt6-\sqrt2}{\sqrt3}
=\frac{50}{3}\left(3\sqrt2-\sqrt6\right).
\]

\gksubpart{（3）}

\Needspace{5\baselineskip}
以极点为原点、极轴为 $x$ 轴正向，代入 $x=r\cos\theta$、$y=r\sin\theta$：

\nopagebreak[4]
\[
r^2\cos^2\theta-r^2\sin^2\theta=a^2,
\qquad\text{\text{即}}\qquad
r^2\cos2\theta=a^2.
\]

\gkanswer{答案（1）：$50$。}

\gkanswer{答案（2）：$c=\dfrac{50}{3}\left(3\sqrt2-\sqrt6\right)$。}

\gkanswer{答案（3）：$r^2\cos2\theta=a^2$。}

\end{answer}
\begin{answer}
\gksolhead{6}{}

\textbf{题目}：

$\triangle ABC$ 的 $\angle A$ 的平分线交 $BC$ 于 $D$，延长 $AD$ 交三角形的外接圆于 $E$，求证：$AB\cdot AC=AD\cdot AE$。

\Needspace{4\baselineskip}
\textbf{解析}

在 $\triangle ABD$ 与 $\triangle ACE$ 中，先证两对角对应相等。

\Needspace{5\baselineskip}
其一，$AD$ 平分 $\angle A$ 且 $A$、$D$、$E$ 共线，故

\nopagebreak[4]
\[
\angle BAD=\angle CAE.
\]

\Needspace{5\baselineskip}
其二，$\angle ADB$ 是 $\triangle ADC$ 在 $D$ 处的外角：

\nopagebreak[4]
\[
\angle ADB=\angle DAC+\angle DCA=\angle DAC+\angle BCA.
\]

\Needspace{5\baselineskip}
又 $A$、$B$、$C$、$E$ 四点共圆，$\angle BCE$ 与 $\angle BAE$ 是同弧 $BE$ 上的圆周角，故 $\angle BCE=\angle BAE=\angle BAD=\angle DAC$，从而

\nopagebreak[4]
\[
\angle ACE=\angle BCA+\angle BCE=\angle BCA+\angle DAC=\angle ADB.
\]

\Needspace{5\baselineskip}
两角对应相等，得

\nopagebreak[4]
\[
\triangle ABD\sim\triangle AEC
\]

\Needspace{5\baselineskip}
（对应关系 $A\leftrightarrow A$、$B\leftrightarrow E$、$D\leftrightarrow C$），于是

\nopagebreak[4]
\[
\frac{AB}{AE}=\frac{AD}{AC},
\qquad\text{\text{即}}\qquad
AB\cdot AC=AD\cdot AE.
\]

\gkanswer{答案：见解析版}

\end{answer}
\begin{answer}
\gksolhead{7}{}

\textbf{题目}：

已知正三棱锥 $P\text{—}ABC$ 的底面边长是 $6\ \mathrm{cm}$，侧面与底面所成的角是 $60^\circ$，求它的体积。（结果可用根式表示）

\Needspace{4\baselineskip}
\textbf{解析}

\Needspace{5\baselineskip}
设顶点 $P$ 在底面内的射影为底面正三角形的中心 $O$，$M$ 为 $AB$ 的中点。底面边长为 $6\ \mathrm{cm}$，则底面积与边心距分别为

\nopagebreak[4]
\[
S=\frac{\sqrt3}{4}\times6^2=9\sqrt3\ \mathrm{cm^2},
\qquad
OM=\frac{\sqrt3}{6}\times6=\sqrt3\ \mathrm{cm}.
\]

\Needspace{5\baselineskip}
侧面与底面所成的二面角的平面角即 $\angle PMO=60^\circ$，故棱锥的高

\nopagebreak[4]
\[
h=PO=OM\cdot\operatorname{tg}60^\circ=\sqrt3\times\sqrt3=3\ \mathrm{cm}.
\]

\Needspace{5\baselineskip}
体积

\nopagebreak[4]
\[
V=\frac13Sh=\frac13\times9\sqrt3\times3=9\sqrt3\ \mathrm{cm^3}.
\]

\gkanswer{答案：$V=9\sqrt3\ \mathrm{cm^3}$。}

\end{answer}
\begin{answer}
\gksolhead{8}{}

\textbf{题目}：

\Needspace{5\baselineskip}
已知 $a$、$b$、$c$ 为三角形的三条边，求证：

\nopagebreak[4]
\[
b^2x^2+(b^2+c^2-a^2)x+c^2=0
\]

无实数根。

\begin{gktip}
\textbf{源注}：原卷二次项印作“$b^2x^2$”，其指数 2 处带一墨点（易误读为下标）；按二次方程题意并结合 600 dpi 局部重渲确认为 $x^2$，墨点系印刷污渍。
\end{gktip}

\Needspace{4\baselineskip}
\textbf{解析}

\Needspace{5\baselineskip}
方程是关于 $x$ 的二次方程（首项系数 $b^2>0$），其判别式为

\nopagebreak[4]
\[
\Delta=(b^2+c^2-a^2)^2-4b^2c^2.
\]

\Needspace{5\baselineskip}
由余弦定理，$b^2+c^2-a^2=2bc\cos A$，代入：

\nopagebreak[4]
\[
\Delta=4b^2c^2\cos^2A-4b^2c^2
=-4b^2c^2\sin^2A.
\]

\Needspace{5\baselineskip}
因 $a$、$b$、$c$ 为三角形的三条边，$b>0$、$c>0$，且 $0<A<\pi$、$\sin A>0$，故

\nopagebreak[4]
\[
\Delta=-4b^2c^2\sin^2A<0,
\]

方程无实数根。

\gkanswer{答案：见解析版}

\end{answer}
\begin{answer}
\gksolhead{9}{}

\textbf{题目}：

方程 $x^2-(\operatorname{tg}\theta+\operatorname{ctg}\theta)x+1=0$ 有一根为 $2+\sqrt3$，求 $\sin2\theta$。

\Needspace{4\baselineskip}
\textbf{解析}

\Needspace{5\baselineskip}
设方程的两根为 $x_1=2+\sqrt3$ 与 $x_2$。由韦达定理，$x_1x_2=1$，故

\nopagebreak[4]
\[
x_2=\frac{1}{2+\sqrt3}=2-\sqrt3.
\]

\Needspace{5\baselineskip}
两根之和为

\nopagebreak[4]
\[
x_1+x_2=(2+\sqrt3)+(2-\sqrt3)=4.
\]

\Needspace{5\baselineskip}
另一方面，

\nopagebreak[4]
\[
\operatorname{tg}\theta+\operatorname{ctg}\theta
=\frac{\sin\theta}{\cos\theta}+\frac{\cos\theta}{\sin\theta}
=\frac{1}{\sin\theta\cos\theta}
=\frac{2}{\sin2\theta}.
\]

\Needspace{5\baselineskip}
由韦达定理，$\operatorname{tg}\theta+\operatorname{ctg}\theta=x_1+x_2=4$，故

\nopagebreak[4]
\[
\frac{2}{\sin2\theta}=4
\quad\Longrightarrow\quad
\sin2\theta=\frac12.
\]

\gkanswer{答案：$\sin2\theta=\dfrac12$。}

\end{answer}
\begin{answer}
\gksolhead{10}{}

\textbf{题目}：

汽车与自行车分别从相距 75 公里的 $A$、$B$ 两地同时出发，（$A$ 地在 $B$ 地的正东方向）汽车从 $A$ 地向正西方向行驶，自行车从 $B$ 地向正南方向行驶，已知汽车和自行车的速度分别为 30 公里/小时和 15 公里/小时，问出发后多少小时两车相距最短？

\Needspace{4\baselineskip}
\textbf{解析}

\Needspace{5\baselineskip}
以 $B$ 为原点、正东方向为 $x$ 轴正向建立直角坐标系，则 $A(75,0)$。设出发 $t$ 小时后，汽车位于 $(75-30t,\ 0)$，自行车位于 $(0,\ -15t)$。两车距离的平方

\nopagebreak[4]
\[
d^2=(75-30t)^2+(15t)^2
=5625-4500t+900t^2+225t^2
=1125t^2-4500t+5625.
\]

\Needspace{5\baselineskip}
配方：

\nopagebreak[4]
\[
d^2=1125(t-2)^2+1125.
\]

\Needspace{5\baselineskip}
故 $t=2$ 时 $d^2$ 取最小值 $1125$，最短距离

\nopagebreak[4]
\[
d=\sqrt{1125}=15\sqrt5\approx33.5\ \text{\text{公里}}.
\]

即出发后 2 小时两车相距最短。

\gkanswer{答案：出发后 $2$ 小时两车相距最短，最短距离为 $15\sqrt5$ 公里（约 $33.5$ 公里）。}

\end{answer}
\begin{answer}
\gksolhead{11}{}

\textbf{题目}：

已知等差数列有 10 项，偶数项之和为 55，奇数项之和为 45，求此数列的第五项。

\Needspace{4\baselineskip}
\textbf{解析}

\Needspace{5\baselineskip}
十项中奇数项 $a_1$、$a_3$、$a_5$、$a_7$、$a_9$ 关于 $a_5$ 对称，偶数项 $a_2$、$a_4$、$a_6$、$a_8$、$a_{10}$ 关于 $a_6$ 对称，由等差数列的性质：

\nopagebreak[4]
\[
a_1+a_3+a_5+a_7+a_9=5a_5=45
\quad\Longrightarrow\quad
a_5=9,
\]

\nopagebreak[4]
\[
a_2+a_4+a_6+a_8+a_{10}=5a_6=55
\quad\Longrightarrow\quad
a_6=11.
\]

（顺带可得公差 $d=a_6-a_5=2$。）所求第五项 $a_5=9$。

\gkanswer{答案：第五项 $a_5=9$。}

\end{answer}
\begin{answer}
\gksolhead{12}{}

\textbf{题目}：

一动点到直线 $x=8$ 的距离是它到点 $A(2,0)$ 的距离的 2 倍，

（1） 求这动点的轨迹方程；

（2） 判断这轨迹是何曲线；

（3） 求经过这曲线上一点 $B(2,3)$ 的切线方程。

\Needspace{4\baselineskip}
\textbf{解析}

\gksubpart{（1）}

\Needspace{5\baselineskip}
设动点坐标为 $(x,y)$。它到直线 $x=8$ 的距离为 $|x-8|$，到点 $A(2,0)$ 的距离为 $\sqrt{(x-2)^2+y^2}$。依题意

\nopagebreak[4]
\[
|x-8|=2\sqrt{(x-2)^2+y^2}.
\]

\Needspace{5\baselineskip}
两边平方：

\nopagebreak[4]
\[
x^2-16x+64=4\left(x^2-4x+4+y^2\right),
\]

\Needspace{5\baselineskip}
整理得

\nopagebreak[4]
\[
3x^2+4y^2-48=0,
\qquad
\frac{x^2}{16}+\frac{y^2}{12}=1.
\]

\gksubpart{（2）}

轨迹方程是中心在原点、焦点在 $x$ 轴上的椭圆：$a^2=16$、$b^2=12$，$c^2=a^2-b^2=4$，焦点为 $(\pm2,0)$。

\gksubpart{（3）}

\Needspace{5\baselineskip}
先验证 $B(2,3)$ 在曲线上：$\dfrac{4}{16}+\dfrac{9}{12}=1$，成立。过椭圆 $\dfrac{x^2}{a^2}+\dfrac{y^2}{b^2}=1$ 上一点 $(x_0,y_0)$ 的切线方程为 $\dfrac{x_0x}{a^2}+\dfrac{y_0y}{b^2}=1$，代入 $B(2,3)$：

\nopagebreak[4]
\[
\frac{2x}{16}+\frac{3y}{12}=1
\quad\Longrightarrow\quad
x+2y-8=0.
\]

（亦可设过 $B$ 的直线为 $y-3=k(x-2)$，与椭圆方程联立后令判别式为零，解得 $k=-\dfrac12$，化简得同一切线。）

\gkanswer{答案（1）：$\dfrac{x^2}{16}+\dfrac{y^2}{12}=1$。}

\gkanswer{答案（2）：椭圆（$a=4$、$b=2\sqrt3$，焦点 $(\pm2,0)$ 在 $x$ 轴上）。}

\gkanswer{答案（3）：$x+2y-8=0$。}

\end{answer}